\documentclass[12pt]{article}

\usepackage[T1]{fontenc}
\usepackage[utf8]{inputenc}
\usepackage[brazil]{babel}
\usepackage{lmodern}
\usepackage{microtype}
\usepackage[a4paper,margin=2.5cm]{geometry}
\usepackage{setspace}
\usepackage{indentfirst}

\setstretch{1.15}

\title{O Papel da Política Monetária \thanks{Texto Traduzido livremente por Luiz Mario Andrade com o auxílio da ferramenta Chat GPT 5.4}}
\author{Milton Friedman}
\date{Março de 1968}

\begin{document}

\maketitle

Há um amplo acordo sobre os principais objetivos da política econômica: alto nível de emprego, preços estáveis e crescimento rápido. Há menos acordo sobre se esses objetivos são mutuamente compatíveis ou, entre aqueles que os consideram incompatíveis, sobre os termos em que eles podem e devem ser substituídos uns pelos outros. Há o menor acordo sobre o papel que os vários instrumentos de política podem e devem desempenhar no alcance desses diversos objetivos.

Meu tópico para esta noite é o papel de um desses instrumentos — a política monetária. O que ela pode contribuir? E como deve ser conduzida para contribuir ao máximo? A opinião sobre essas questões tem flutuado amplamente. No primeiro surto de entusiasmo sobre o recém-criado Sistema de Reserva Federal, muitos observadores atribuíram a estabilidade relativa da década de 1920 à capacidade do Sistema para o "ajuste fino" (*fine tuning*) — para aplicar um termo moderno apto. Passou-se a acreditar amplamente que uma nova era havia chegado, na qual os ciclos econômicos haviam se tornado obsoletos devido aos avanços na tecnologia monetária. Essa opinião era compartilhada por economistas e leigos, embora, é claro, houvesse algumas vozes dissonantes. A Grande Contração destruiu essa atitude ingênua. A opinião oscilou para o outro extremo. A política monetária era um barbante. Você podia puxá-lo para interromper a inflação, mas não podia empurrá-lo para deter a recessão. Você poderia levar um cavalo até a água, mas não poderia obrigá-lo a beber. Tal teoria por aforismo logo foi substituída pela análise rigorosa e sofisticada de Keynes.

Keynes ofereceu simultaneamente uma explicação para a presumida impotência da política monetária para conter a depressão, uma interpretação não monetária da depressão e uma alternativa à política monetária para enfrentar a depressão, e sua oferta foi avidamente aceita. Se a preferência pela liquidez é absoluta ou quase absoluta — como Keynes acreditava provável em tempos de desemprego pesado — as taxas de juros não podem ser reduzidas por medidas monetárias. Se o investimento e o consumo são pouco afetados pelas taxas de juros — como Hansen e muitos outros discípulos americanos de Keynes passaram a acreditar — taxas de juros mais baixas, mesmo que pudessem ser alcançadas, fariam pouco bem. A política monetária é duplamente condenada. A contração, desencadeada, sob essa visão, por um colapso do investimento ou por uma escassez de oportunidades de investimento ou por uma parcimônia obstinada, não poderia, argumentava-se, ter sido detida por medidas monetárias. Mas havia uma alternativa disponível — a política fiscal. Os gastos do governo poderiam compensar o investimento privado insuficiente. Reduções de impostos poderiam minar a parcimônia obstinada.

A ampla aceitação dessas visões na profissão econômica significou que, por cerca de duas décadas, a política monetária foi considerada por todos, exceto por algumas almas reacionárias, como tendo sido tornada obsoleta pelo novo conhecimento econômico. O dinheiro não importava. Seu único papel era o papel menor de manter as taxas de juros baixas, a fim de conter os pagamentos de juros no orçamento do governo, contribuir para a "eutanásia do rentista" e, talvez, estimular um pouco o investimento para auxiliar os gastos do governo na manutenção de um alto nível de demanda agregada.

Essas visões produziram uma adoção generalizada de políticas de "dinheiro barato" após a guerra. E elas receberam um choque rude quando essas políticas falharam em país após país, quando banco central após banco central foi forçado a desistir da pretensão de que poderia manter indefinidamente "a" taxa de juros em um nível baixo. Neste país, o desfecho público veio com o Acordo Federal Reserve-Tesouro em 1951, embora a política de fixar os preços dos títulos do governo não tenha sido formalmente abandonada até 1953. A inflação, estimulada por políticas de dinheiro barato, e não a amplamente anunciada depressão do pós-guerra, acabou sendo a ordem do dia. O resultado foi o início de um renascimento da crença na potência da política monetária.

Este renascimento foi fortemente incentivado entre os economistas pelos desenvolvimentos teóricos iniciados por Haberler, mas nomeados em homenagem a Pigou, que apontaram um canal — a saber, mudanças na riqueza — por meio do qual mudanças na quantidade real de dinheiro podem afetar a demanda agregada, mesmo que não alterem as taxas de juros. Esses desenvolvimentos teóricos não minaram o argumento de Keynes contra a potência das medidas monetárias ortodoxas quando a preferência pela liquidez é absoluta, uma vez que, sob tais circunstâncias, as operações monetárias usuais envolvem simplesmente a substituição de dinheiro por outros ativos sem alterar a riqueza total. Mas eles mostraram como mudanças na quantidade de dinheiro produzidas de outras formas poderiam afetar os gastos totais mesmo sob tais circunstâncias. E, mais fundamentalmente, eles minaram a proposição teórica central de Keynes, a saber, que mesmo em um mundo de preços flexíveis, uma posição de equilíbrio em pleno emprego poderia não existir. Doravante, o desemprego teve novamente que ser explicado por rigidezes ou imperfeições, não como o resultado natural de um processo de mercado plenamente operante.

O renascimento da crença na potência da política monetária foi fomentado também por uma reavaliação do papel que o dinheiro desempenhou de 1929 a 1933. Keynes e a maioria dos outros economistas da época acreditavam que a Grande Contração nos Estados Unidos ocorreu apesar de políticas expansionistas agressivas por parte das autoridades monetárias — que elas fizeram o seu melhor, mas o seu melhor não foi bom o suficiente. Estudos recentes demonstraram que os fatos são precisamente o oposto: as autoridades monetárias dos EUA seguiram políticas altamente deflacionárias. A quantidade de dinheiro nos Estados Unidos caiu um terço no curso da contração. E caiu não porque não houvesse tomadores de empréstimos dispostos — não porque o cavalo não bebesse a água. Caiu porque o Sistema de Reserva Federal forçou ou permitiu uma redução acentuada na base monetária, porque falhou em exercer as responsabilidades atribuídas a ele na Lei de Reserva Federal de fornecer liquidez ao sistema bancário. A Grande Contração é um testemunho trágico do poder da política monetária — não, como Keynes e tantos de seus contemporâneos acreditaram, evidência de sua impotência.

Nos Estados Unidos, o renascimento da crença na potência da política monetária foi fortalecido também pela crescente desilusão com a política fiscal, não tanto com seu potencial para afetar a demanda agregada, mas com a viabilidade prática e política de usá-la dessa forma. As despesas revelaram-se responder de forma lenta e com longos atrasos às tentativas de ajustá-las ao curso da atividade econômica, de modo que a ênfase mudou para os impostos. Mas aqui fatores políticos entraram com vingança para impedir o ajuste imediato à necessidade presumida, como foi tão graficamente ilustrado nos meses desde que escrevi o primeiro rascunho desta palestra. "Ajuste fino" é uma frase maravilhosamente evocativa nesta era eletrônica, mas tem pouca semelhança com o que é possível na prática — o que, eu acrescentaria, não é um mal absoluto.

É difícil perceber quão radical foi a mudança na opinião profissional sobre o papel do dinheiro. Quase nenhum economista hoje aceita visões que eram moeda comum há cerca de duas décadas. Permitam-me citar alguns exemplos.

Em uma palestra publicada em 1945, E. A. Goldenweiser, então Diretor da Divisão de Pesquisa do Conselho do Federal Reserve, descreveu o objetivo principal da política monetária como sendo "manter o valor dos títulos do Governo... Este país", escreveu ele, "terá que se ajustar a uma taxa de juros de 2,5 por cento como o retorno sobre o dinheiro seguro e de longo prazo, porque chegou o tempo em que os retornos sobre o capital pioneiro não podem mais ser ilimitados como eram no passado" [4, p. 117].

Em um livro sobre *Financing American Prosperity*, editado por Paul Homan e Fritz Machlup e publicado em 1945, Alvin Hansen dedica nove páginas de texto ao "problema poupança-investimento" sem encontrar qualquer necessidade de usar as palavras "taxa de juros" ou qualquer fac-símile próximo a elas [5, pp. 218-27]. Em sua contribuição para este volume, Fritz Machlup escreveu: "Questões relativas à taxa de juros, em particular em relação à sua variação ou à sua estabilidade, podem não estar entre os problemas mais vitais da economia do pós-guerra, mas são certamente algumas das mais perplexas" [5, p. 466]. Em sua contribuição, John H. Williams — não apenas professor em Harvard, mas também consultor de longa data do Federal Reserve Bank de Nova York — escreveu: "Não vejo perspectiva de renascimento de um controle monetário geral no período do pós-guerra" [5, p. 383].

Outro dos volumes que tratam da política do pós-guerra que apareceu nesta época, *Planning and Paying for Full Employment*, foi editado por Abba P. Lerner e Frank D. Graham [6] e contou com colaboradores de todos os matizes da opinião profissional — de Henry Simons e Frank Graham a Abba Lerner e Hans Neisser. No entanto, Albert Halasi, em seu excelente resumo dos artigos, pôde dizer: "Nossos colaboradores não discutem a questão da oferta de moeda... Os colaboradores não fazem menção especial à política de crédito para remediar depressões reais... A inflação... poderia ser combatida de forma mais eficaz aumentando as taxas de juros... Mas... outras medidas anti-inflacionárias... são preferíveis" [6, pp. 23-24]. *A Survey of Contemporary Economics*, editado por Howard Ellis e publicado em 1948, foi uma tentativa "oficial" de codificar o estado do pensamento econômico da época. Em sua contribuição, Arthur Smithies escreveu: "No campo da ação compensatória, acredito que a política fiscal deve assumir a maior parte da carga. Sua principal rival, a política monetária, parece estar desqualificada por motivos institucionais. Este país parece estar comprometido com algo como o atual nível baixo de taxas de juros em uma base de longo prazo" [1, p. 208].

Essas citações sugerem o sabor do pensamento profissional de cerca de duas décadas atrás. Se você deseja ir mais longe nesta investigação humilhante, recomendo que compare as seções sobre dinheiro — quando conseguir encontrá-las — nos textos de Princípios dos primeiros anos do pós-guerra com as longas seções na safra atual, mesmo, ou especialmente, quando os Princípios iniciais e recentes são edições diferentes da mesma obra.

O pêndulo oscilou muito desde então, se não totalmente para a posição do final da década de 1920, pelo menos muito mais perto dessa posição do que da posição de 1945. Há, é claro, muitas diferenças entre então e agora, menos na potência atribuída à política monetária do que nos papéis atribuídos a ela e nos critérios pelos quais a profissão acredita que a política monetária deva ser guiada. Naquela época, os principais papéis atribuídos à política monetária eram promover a estabilidade de preços e preservar o padrão-ouro; os principais critérios da política monetária eram o estado do "mercado monetário", a extensão da "especulação" e o movimento do ouro. Hoje, a primazia é atribuída à promoção do pleno emprego, sendo a prevenção da inflação um objetivo contínuo, mas definitivamente secundário. E há um grande desacordo sobre os critérios de política, variando da ênfase nas condições do mercado monetário, taxas de juros e quantidade de dinheiro à crença de que o próprio estado do emprego deve ser o critério imediato da política.

Ressalto, no entanto, a semelhança entre as visões que prevaleciam no final dos anos 20 e as que prevalecem hoje porque temo que, agora como então, o pêndulo possa muito bem ter oscilado demais; que, agora como então, estejamos em perigo de atribuir à política monetária um papel maior do que ela pode desempenhar, em perigo de pedir que ela realize tarefas que não pode alcançar e, como resultado, em perigo de impedi-la de fazer a contribuição que é capaz de fazer.

Desacostumado como estou a denegrir a importância do dinheiro, portanto, terei como minha primeira tarefa enfatizar o que a política monetária não pode fazer. Tentarei então delinear o que ela pode fazer e como pode melhor fazer sua contribuição, no estado atual do nosso conhecimento — ou ignorância.

\section*{I. O Que a Política Monetária Não Pode Fazer}

Do mundo infinito da negação, selecionei duas limitações da política monetária para discutir: (1) Ela não pode fixar as taxas de juros por mais de períodos muito limitados; (2) Ela não pode fixar a taxa de desemprego por mais de períodos muito limitados. Seleciono estas porque o contrário tem sido ou é amplamente acreditado, porque correspondem às duas principais tarefas inalcançáveis que são muito prováveis de serem atribuídas à política monetária, e porque essencialmente a mesma análise teórica cobre ambas.

\subsection*{Fixação das Taxas de Juros}

A história já persuadiu muitos de vocês sobre a primeira limitação. Como observado anteriormente, o fracasso das políticas de dinheiro barato foi uma fonte importante da reação contra o keynesianismo simplista. Nos Estados Unidos, essa reação envolveu o reconhecimento generalizado de que a fixação dos preços dos títulos durante e após a guerra foi um erro, que o abandono dessa política foi um passo desejável e inevitável, e que não teve nenhuma das consequências perturbadoras e desastrosas que foram tão livremente previstas na época.

A limitação deriva de uma característica muito mal compreendida da relação entre dinheiro e taxas de juros. Deixe o Fed se propor a manter as taxas de juros baixas. Como ele tentará fazer isso? Comprando títulos. Isso aumenta seus preços e reduz seus rendimentos. No processo, também aumenta a quantidade de reservas disponíveis para os bancos, consequentemente o montante de crédito bancário e, por fim, a quantidade total de dinheiro. É por isso que os banqueiros centrais em particular, e a comunidade financeira em geral, acreditam que um aumento na quantidade de dinheiro tende a reduzir as taxas de juros. Os economistas acadêmicos aceitam a mesma conclusão, mas por razões diferentes. Eles veem, em sua mente, um cronograma de preferência pela liquidez com inclinação negativa. Como as pessoas podem ser induzidas a manter uma maior quantidade de dinheiro? Apenas pela redução das taxas de juros.

Ambos estão certos, até certo ponto. O impacto inicial do aumento da quantidade de dinheiro a uma taxa mais rápida do que tem aumentado é tornar as taxas de juros mais baixas por um tempo do que seriam de outra forma. Mas este é apenas o começo do processo, não o fim. A taxa mais rápida de crescimento monetário estimulará os gastos, tanto através do impacto no investimento de taxas de juros de mercado mais baixas quanto através do impacto em outros gastos e, consequentemente, nos preços relativos de saldos de caixa mais elevados do que o desejado. Mas o gasto de um homem é a renda de outro homem. A renda crescente elevará o cronograma de preferência pela liquidez e a demanda por empréstimos; também pode elevar os preços, o que reduziria a quantidade real de dinheiro. Esses três efeitos reverterão a pressão inicial de queda nas taxas de juros de forma bastante rápida, digamos, em algo menos de um ano. Juntos, eles tenderão, após um intervalo um pouco mais longo, digamos, um ou dois anos, a retornar as taxas de juros ao nível que teriam de outra forma. De fato, dada a tendência da economia de reagir exageradamente, é altamente provável que elevem as taxas de juros temporariamente além desse nível, desencadeando um processo de ajuste cíclico.

Um quarto efeito, se e quando se tornar operativo, irá ainda mais longe e significará definitivamente que uma taxa mais alta de expansão monetária corresponderá a um nível de taxas de juros mais alto, e não mais baixo, do que teria prevalecido de outra forma. Deixe a taxa mais alta de crescimento monetário produzir preços crescentes e deixe o público passar a esperar que os preços continuem a subir. Os tomadores de empréstimos estarão então dispostos a pagar e os credores exigirão taxas de juros mais altas — como Irving Fisher apontou décadas atrás. Esse efeito de expectativa de preço é lento para se desenvolver e também lento para desaparecer. Fisher estimou que levava várias décadas para um ajuste completo e trabalhos mais recentes são consistentes com suas estimativas.

Esses efeitos subsequentes explicam por que toda tentativa de manter as taxas de juros em um nível baixo forçou a autoridade monetária a se envolver em compras sucessivamente maiores e maiores no mercado aberto. Eles explicam por que, historicamente, taxas de juros nominais altas e crescentes têm sido associadas ao rápido crescimento na quantidade de dinheiro, como no Brasil ou no Chile ou nos Estados Unidos nos últimos anos, e por que taxas de juros baixas e cadentes têm sido associadas ao crescimento lento na quantidade de dinheiro, como na Suíça agora ou nos Estados Unidos de 1929 a 1933. Como uma questão empírica, taxas de juros baixas são um sinal de que a política monetária tem sido \textit{apertada} — no sentido de que a quantidade de dinheiro cresceu lentamente; taxas de juros altas são um sinal de que a política monetária tem sido \textit{frouxa} — no sentido de que a quantidade de dinheiro cresceu rapidamente. Os fatos mais amplos da experiência correm precisamente na direção oposta àquela que a comunidade financeira e os economistas acadêmicos geralmente consideram como certa.

Paradoxalmente, a autoridade monetária poderia garantir taxas nominais de juros baixas — mas para fazer isso teria que começar no que parece ser a direção oposta, engajando-se em uma política monetária deflacionária. Da mesma forma, poderia garantir taxas de juros nominais altas engajando-se em uma política inflacionária e aceitando um movimento temporário nas taxas de juros na direção oposta.

Essas considerações não apenas explicam por que a política monetária não pode fixar as taxas de juros; elas também explicam por que as taxas de juros são um indicador tão enganoso de se a política monetária está "apertada" ou "frouxa". Para isso, é muito melhor olhar para a taxa de variação da quantidade de dinheiro.

\subsection*{O Emprego como Critério de Política}

A segunda limitação que desejo discutir vai mais contra a corrente do pensamento atual. O crescimento monetário, acredita-se amplamente, tenderá a estimular o emprego; a contração monetária, a retardar o emprego. Por que, então, a autoridade monetária não pode adotar uma meta para o emprego ou desemprego — digamos, 3 por cento de desemprego; ser apertada quando o desemprego for menor que a meta; ser frouxa quando o desemprego for maior que a meta; e desta forma fixar o desemprego em, digamos, 3 por cento? A razão pela qual não pode é precisamente a mesma que para as taxas de juros — a diferença entre as consequências imediatas e as retardadas de tal política.

Graças a Wicksell, todos estamos familiarizados com o conceito de uma taxa de juros "natural" e a possibilidade de uma discrepância entre a taxa "natural" e a de "mercado". A análise precedente das taxas de juros pode ser traduzida diretamente em termos wicksellianos. A autoridade monetária pode tornar a taxa de mercado menor que a taxa natural apenas por meio da inflação. Pode tornar a taxa de mercado mais alta que a taxa natural apenas por meio da deflação. Adicionamos apenas uma ruga a Wicksell — a distinção de Irving Fisher entre a taxa de juros nominal e a real. Deixe a autoridade monetária manter a taxa nominal de mercado por um tempo abaixo da taxa natural por meio da inflação. Isso, por sua vez, elevará a própria taxa natural nominal, uma vez que as antecipações de inflação se tornem generalizadas, exigindo assim uma inflação ainda mais rápida para conter a taxa de mercado. Da mesma forma, por causa do efeito Fisher, exigirá não apenas deflação, mas uma deflação cada vez mais rápida para manter a taxa de mercado acima da taxa "natural" inicial.

Esta análise tem sua contraparte próxima no mercado de trabalho. Em qualquer momento do tempo, há algum nível de desemprego que tem a propriedade de ser consistente com o equilíbrio na estrutura das taxas de salários \textit{reais}. Nesse nível de desemprego, as taxas de salários reais tendem, em média, a subir a uma taxa secular "normal", ou seja, a uma taxa que pode ser mantida indefinidamente enquanto a formação de capital, as melhorias tecnológicas, etc., permanecerem em suas tendências de longo prazo. Um nível mais baixo de desemprego é uma indicação de que há um excesso de demanda por trabalho que produzirá pressão ascendente sobre as taxas de salários reais. Um nível mais alto de desemprego é uma indicação de que há um excesso de oferta de trabalho que produzirá pressão descendente sobre as taxas de salários reais. A "taxa natural de desemprego", em outras palavras, é o nível que seria gerado pelo sistema walrasiano de equações de equilíbrio geral, desde que nele estejam inseridas as características estruturais reais dos mercados de trabalho e de commodities, incluindo imperfeições de mercado, variabilidade estocástica em demandas e ofertas, o custo de reunir informações sobre vagas de emprego e disponibilidades de mão de obra, os custos de mobilidade, e assim por diante.

Vocês reconhecerão a estreita semelhança entre esta afirmação e a célebre Curva de Phillips. A semelhança não é coincidência. A análise de Phillips sobre a relação entre desemprego e mudança salarial é merecidamente celebrada como uma contribuição importante e original. Mas, infelizmente, ela contém um defeito básico — a falha em distinguir entre salários nominais e salários reais — assim como a análise de Wicksell falhou em distinguir entre taxas de juros nominais e taxas de juros reais. Implicitamente, Phillips escreveu seu artigo para um mundo no qual todos antecipavam que os preços nominais seriam estáveis e no qual essa antecipação permanecia inabalável e imutável, independentemente do que acontecesse com os preços e salários reais. Suponha, ao contrário, que todos antecipem que os preços subirão a uma taxa de mais de 75 por cento ao ano — como, por exemplo, os brasileiros faziam há alguns anos. Então os salários devem subir a essa taxa simplesmente para manter os salários reais inalterados. Um excesso de oferta de mão de obra refletir-se-á em um aumento menos rápido nos salários nominais do que nos preços antecipados, e não em um declínio absoluto nos salários. Quando o Brasil embarcou em uma política para reduzir a taxa de aumento de preços e conseguiu baixar o aumento de preços para cerca de 45 por cento ao ano, houve um aumento inicial acentuado no desemprego porque, sob a influência de antecipações anteriores, os salários continuaram subindo a um ritmo que era superior à nova taxa de aumento de preços, embora inferior à anterior. Este é o resultado experimentado, e esperado, de todas as tentativas de reduzir a taxa de inflação abaixo daquela amplamente antecipada.

Para evitar mal-entendidos, permitam-me enfatizar que, ao usar o termo taxa "natural" de desemprego, não pretendo sugerir que ela seja imutável e inalterável. Pelo contrário, muitas das características do mercado que determinam o seu nível são feitas pelo homem e por políticas. Nos Estados Unidos, por exemplo, as taxas de salário mínimo legal, as leis Walsh-Healy e Davis-Bacon e a força dos sindicatos contribuem para tornar a taxa natural de desemprego mais alta do que seria de outra forma. Melhorias nas agências de emprego, na disponibilidade de informações sobre vagas de emprego e oferta de mão de obra, e assim por diante, tenderiam a baixar a taxa natural de desemprego. Uso o termo "natural" pela mesma razão que Wicksell o fez — para tentar separar as forças reais das forças monetárias.

Vamos supor que a autoridade monetária tente fixar a taxa de desemprego "de mercado" em um nível abaixo da taxa "natural". Para maior clareza, suponha que ela tome 3 por cento como a taxa alvo e que a taxa "natural" seja superior a 3 por cento. Suponha também que comecemos em um momento em que os preços estiveram estáveis e quando o desemprego seja superior a 3 por cento. Consequentemente, a autoridade aumenta a taxa de crescimento monetário. Isso será expansionista. Ao tornar os saldos nominais de caixa superiores aos desejados pelas pessoas, tenderá inicialmente a reduzir as taxas de juros e, desta e de outras formas, a estimular os gastos. A renda e os gastos começarão a subir.

Para começar, grande parte ou a maior parte do aumento na renda assumirá a forma de um aumento na produção e no emprego, em vez de nos preços. As pessoas têm esperado que os preços sejam estáveis, e os preços e salários foram estabelecidos por algum tempo no futuro com base nisso. Leva tempo para as pessoas se ajustarem a um novo estado de demanda. Os produtores tenderão a reagir à expansão inicial na demanda agregada aumentando a produção, os funcionários trabalhando mais horas e os desempregados aceitando empregos agora oferecidos aos salários nominais anteriores. Isso é doutrina bastante padrão.

Mas descreve apenas os efeitos iniciais. Como os preços de venda dos produtos tipicamente respondem a um aumento não antecipado na demanda nominal mais rápido do que os preços dos fatores de produção, os salários reais recebidos caíram — embora os salários reais antecipados pelos funcionários tenham subido, uma vez que os funcionários avaliaram implicitamente os salários oferecidos ao nível de preços anterior. De fato, a queda simultânea \textit{ex post} nos salários reais para os empregadores e o aumento \textit{ex ante} nos salários reais para os funcionários é o que permitiu o aumento do emprego. Mas o declínio \textit{ex post} nos salários reais logo passará a afetar as antecipações. Os funcionários começarão a contar com preços crescentes das coisas que compram e a exigir salários nominais mais altos para o futuro. O desemprego "de mercado" está abaixo do nível "natural". Há um excesso de demanda por trabalho, de modo que os salários reais tenderão a subir em direção ao seu nível inicial.

Mesmo que a taxa mais elevada de crescimento monetário continue, o aumento dos salários reais reverterá o declínio no desemprego e, em seguida, levará a um aumento, que tenderá a retornar o desemprego ao seu nível anterior. Para manter o desemprego em seu nível alvo de 3 por cento, a autoridade monetária teria que aumentar ainda mais o crescimento monetário. Como no caso da taxa de juros, a taxa "de mercado" pode ser mantida abaixo da taxa "natural" apenas por meio da inflação. E, como no caso da taxa de juros também, apenas acelerando a inflação. Por outro lado, deixe a autoridade monetária escolher uma taxa alvo de desemprego que esteja acima da taxa natural, e ela será levada a produzir uma deflação, e uma deflação acelerada nisso.

E se a autoridade monetária escolhesse a taxa "natural" — seja de juros ou de desemprego — como seu alvo? Um problema é que ela não pode saber qual é a taxa "natural". Infelizmente, ainda não concebemos nenhum método para estimar com precisão e rapidez a taxa natural de juros ou de desemprego. E a taxa "natural" mudará de tempos em tempos. Mas o problema básico é que mesmo que a autoridade monetária soubesse a taxa "natural" e tentasse fixar a taxa de mercado nesse nível, ela não seria levada a uma política determinada. A taxa "de mercado" variará da taxa natural por todos os tipos de razões que não a política monetária. Se a autoridade monetária responder a essas variações, ela desencadeará efeitos de longo prazo que tornarão qualquer caminho de crescimento monetário que ela siga ultimamente consistente com a regra da política. O curso real do crescimento monetário será análogo a um passeio aleatório (*random walk*), fustigado de um lado para o outro pelas forças que produzem desvios temporários da taxa de mercado em relação à taxa natural.

Para declarar esta conclusão de forma diferente, há sempre um *trade-off* temporário entre inflação e desemprego; não há um *trade-off* permanente. O *trade-off* temporário não vem da inflação per se, mas da inflação não antecipada, o que geralmente significa de uma taxa de inflação crescente. A crença generalizada de que existe um *trade-off* permanente é uma versão sofisticada da confusão entre "alto" e "crescente" que todos reconhecemos em formas mais simples. Uma taxa de inflação crescente pode reduzir o desemprego, uma taxa alta não o fará.

Mas quanto tempo, vocês dirão, é "temporário"? Para as taxas de juros, temos algumas evidências sistemáticas sobre quanto tempo cada um dos vários efeitos leva para se processar. Para o desemprego, não temos. Posso, no máximo, arriscar um julgamento pessoal, baseado em algum exame das evidências históricas, de que os efeitos iniciais de uma taxa de inflação mais alta e não antecipada duram algo entre dois a cinco anos; que esse efeito inicial começa a ser revertido; e que um ajuste completo à nova taxa de inflação leva quase tanto tempo para o emprego quanto para as taxas de juros, digamos, um par de décadas. Para as taxas de juros e para o emprego, permitam-me acrescentar uma ressalva. Essas estimativas são para mudanças na taxa de inflação da ordem de magnitude que tem sido experimentada nos Estados Unidos. Para mudanças muito mais substanciais, como as experimentadas em países da América do Sul, todo o processo de ajuste é grandemente acelerado.

Para declarar a conclusão geral ainda de outra forma, a autoridade monetária controla grandezas nominais — diretamente, a quantidade de suas próprias obrigações. Em princípio, ela pode usar esse controle para fixar uma grandeza nominal — uma taxa de câmbio, o nível de preços, o nível nominal da renda nacional, a quantidade de dinheiro por uma ou outra definição — ou para fixar a taxa de variação de uma grandeza nominal — a taxa de inflação ou deflação, a taxa de crescimento ou declínio da renda nacional nominal, a taxa de crescimento da quantidade de dinheiro. Ela não pode usar seu controle sobre as grandezas nominais para fixar uma grandeza real — a taxa de juros real, a taxa de desemprego, o nível da renda nacional real, a quantidade real de dinheiro, a taxa de crescimento da renda nacional real ou a taxa de crescimento da quantidade real de dinheiro.

\section*{II. O Que a Política Monetária Pode Fazer}

A política monetária não pode fixar essas grandezas reais em níveis predeterminados. Mas a política monetária tem e exerce efeitos importantes sobre essas grandezas reais. Uma coisa não é de forma alguma inconsistente com a outra.

Meus próprios estudos de história monetária tornaram-me extremamente simpático ao comentário frequentemente citado, muito revivido e tão amplamente mal compreendido de John Stuart Mill. "Não pode haver...", escreveu ele, "intrinsecamente uma coisa mais insignificante, na economia da sociedade, do que o dinheiro; exceto no caráter de um dispositivo para poupar tempo e trabalho. É uma máquina para fazer de forma rápida e cômoda o que seria feito, embora menos rapidamente e com menos comodidade, sem ela: e, como muitos outros tipos de maquinaria, só exerce uma influência distinta e independente quando sai de ordem" [7, p. 488].

É verdade que o dinheiro é apenas uma máquina, mas é uma máquina extraordinariamente eficiente. Sem ele, não poderíamos ter começado a atingir o crescimento espantoso na produção e no padrão de vida que experimentamos nos últimos dois séculos — mais do que poderíamos ter feito sem aquelas outras máquinas maravilhosas que pontilham nossa paisagem e nos permitem, em grande parte, simplesmente fazer com mais eficiência o que poderia ser feito sem elas a um custo muito maior em trabalho.

Mas o dinheiro tem uma característica que essas outras máquinas não compartilham. Por ser tão onipresente, quando sai de ordem, lança uma chave de fenda na operação de todas as outras máquinas. A Grande Contração é o exemplo mais dramático, mas não o único. Todas as outras contrações importantes neste país foram ou produzidas por desordem monetária ou grandemente exacerbadas por desordem monetária. Toda inflação importante foi produzida pela expansão monetária — principalmente para atender às demandas prementes de guerra que forçaram a criação de dinheiro para suplementar a tributação explícita.

A primeira e mais importante lição que a história ensina sobre o que a política monetária pode fazer — e é uma lição da mais profunda importância — é que a política monetária pode evitar que o próprio dinheiro seja uma fonte principal de perturbação econômica. Isso soa como uma proposição negativa: evitar grandes erros. Em parte, é. A Grande Contração pode não ter ocorrido de forma alguma, e se tivesse ocorrido, teria sido muito menos severa, se a autoridade monetária tivesse evitado erros, ou se os arranjos monetários fossem os de uma época anterior, quando não havia uma autoridade central com o poder de cometer os tipos de erros que o Sistema de Reserva Federal cometeu. Os últimos anos, para chegar mais perto de casa, teriam sido mais estáveis e mais produtivos de bem-estar econômico se o Federal Reserve tivesse evitado mudanças drásticas e erráticas de direção, primeiro expandindo a oferta de moeda a um ritmo indevidamente rápido, depois, no início de 1966, pisando no freio com muita força, depois, no final de 1966, revertendo-se e retomando a expansão até pelo menos novembro de 1967, a um ritmo mais rápido do que pode ser mantido por muito tempo sem uma inflação apreciável.

Mesmo que a proposição de que a política monetária pode evitar que o próprio dinheiro seja uma fonte principal de perturbação econômica fosse uma proposição totalmente negativa, ela não seria por isso menos importante. Acontece, porém, que não é uma proposição totalmente negativa. A máquina monetária tem saído de ordem mesmo quando não houve uma autoridade central com nada parecido com o poder agora possuído pelo Fed. Nos Estados Unidos, o episódio de 1907 e os pânicos bancários anteriores são exemplos de como a máquina monetária pode sair de ordem em grande parte por conta própria. Existe, portanto, uma tarefa positiva e importante para a autoridade monetária — sugerir melhorias na máquina que reduzam as chances de ela sair de ordem, e usar seus próprios poderes para manter a máquina em boa ordem de funcionamento.

Uma segunda coisa que a política monetária pode fazer é fornecer um pano de fundo estável para a economia — manter a máquina bem lubrificada, para continuar a analogia de Mill. O cumprimento da primeira tarefa contribuirá para este objetivo, mas há mais do que isso. Nosso sistema econômico funcionará melhor quando produtores e consumidores, empregadores e empregados, puderem prosseguir com plena confiança de que o nível médio de preços se comportará de uma forma conhecida no futuro — preferencialmente de que será altamente estável. Sob quaisquer arranjos institucionais concebíveis, e certamente sob aqueles que agora prevalecem nos Estados Unidos, há apenas uma quantidade limitada de flexibilidade nos preços e salários. Precisamos conservar essa flexibilidade para alcançar mudanças nos preços e salários relativos que são necessárias para nos ajustarmos às mudanças dinâmicas nos gostos e na tecnologia. Não devemos dissipá-la simplesmente para alcançar mudanças no nível absoluto de preços que não servem a nenhuma função econômica.

Em uma era anterior, o padrão-ouro era confiado para fornecer confiança na estabilidade monetária futura. Em seu apogeu, desempenhou essa função razoavelmente bem. Claramente não o faz mais, uma vez que mal existe um país no mundo que esteja preparado para deixar o padrão-ouro reinar sem controle — e há razões persuasivas pelas quais os países não devem fazê-lo. A autoridade monetária poderia operar como um substituto para o padrão-ouro, se fixasse as taxas de câmbio e o fizesse exclusivamente alterando a quantidade de dinheiro em resposta aos fluxos de balanço de pagamentos sem "esterilizar" superávits ou déficits e sem recorrer a controle cambial aberto ou oculto ou a mudanças em tarifas e cotas. Mas, novamente, embora muitos banqueiros centrais falem dessa maneira, poucos estão de fato dispostos a seguir esse curso — e, novamente, há razões persuasivas pelas quais não devem fazê-lo. Tal política submeteria cada país às vicissitudes não de um padrão-ouro impessoal e automático, mas das políticas — deliberadas ou acidentais — de outras autoridades monetárias.

No mundo de hoje, se a política monetária deve fornecer um pano de fundo estável para a economia, deve fazê-lo empregando deliberadamente seus poderes para esse fim. Voltarei mais tarde a como ela pode fazer isso.

Finalmente, a política monetária pode contribuir para compensar as principais perturbações no sistema econômico decorrentes de outras fontes. Se houver uma exultação secular independente — como a expansão do pós-guerra foi descrita pelos proponentes da estagnação secular — a política monetária pode, em princípio, ajudar a mantê-la sob controle por uma taxa de crescimento monetário mais lenta do que seria desejável de outra forma. Se, como agora, um orçamento federal explosivo ameaça déficits sem precedentes, a política monetária pode manter quaisquer perigos inflacionários sob controle por uma taxa de crescimento monetário mais lenta do que seria desejável de outra forma. Isso significará temporariamente taxas de juros mais altas do que prevaleceriam de outra forma — para permitir que o governo tome emprestado as somas necessárias para financiar o déficit — mas, ao evitar a aceleração da inflação, pode muito bem significar preços mais baixos e taxas de juros nominais mais baixas a longo prazo. Se o fim de uma guerra substancial oferece ao país uma oportunidade de transferir recursos da produção de guerra para a produção de paz, a política monetária pode facilitar a transição por uma taxa de crescimento monetário mais elevada do que seria desejável de outra forma — embora a experiência não seja muito encorajadora de que ela possa fazer isso sem ir longe demais.

Coloquei este ponto por último e o declarei em termos qualificados — como referindo-se a grandes perturbações — porque acredito que a potencialidade da política monetária em compensar outras forças que provocam instabilidade é muito mais limitada do que se acredita comumente. Simplesmente não sabemos o suficiente para sermos capazes de reconhecer perturbações menores quando ocorrem ou para sermos capazes de prever com precisão quais serão os seus efeitos ou qual política monetária é necessária para compensar os seus efeitos. Não sabemos o suficiente para sermos capazes de atingir objetivos declarados por mudanças delicadas, ou mesmo grosseiras, no mix de política monetária e fiscal. Nesta área, particularmente, o ótimo é provável que seja o inimigo do bom. A experiência sugere que o caminho da sabedoria é usar a política monetária explicitamente para compensar outras perturbações apenas quando estas oferecerem um "perigo claro e presente".

\section*{III. Como a Política Monetária Deve Ser Conduzida?}

Como a política monetária deve ser conduzida para dar a contribuição que é capaz de dar aos nossos objetivos? Esta não é claramente a ocasião para apresentar um "Programa para a Estabilidade Monetária" detalhado — para usar o título de um livro no qual tentei fazer isso [3]. Vou restringir-me aqui a dois requisitos principais para a política monetária que decorrem de forma bastante direta da discussão anterior.

O primeiro requisito é que a autoridade monetária deve guiar-se por grandezas que possa controlar, não por aquelas que não pode controlar. Se, como a autoridade frequentemente fez, ela toma as taxas de juros ou a porcentagem atual de desemprego como o critério imediato da política, será como um veículo espacial que fixou o curso em uma estrela errada. Não importa quão sensível e sofisticado seja o seu aparelho de orientação, o veículo espacial se perderá. E o mesmo acontecerá com a autoridade monetária. Das várias grandezas alternativas que ela pode controlar, os guias mais atraentes para a política são as taxas de câmbio, o nível de preços conforme definido por algum índice e a quantidade de um total monetário — moeda mais depósitos à vista ajustados, ou este total mais depósitos a prazo em bancos comerciais, ou um total ainda mais amplo.

Para os Estados Unidos em particular, as taxas de câmbio são um guia indesejável. Pode valer a pena exigir que a maior parte da economia se ajuste à pequena porcentagem que consiste no comércio exterior se isso garantisse a liberdade da irresponsabilidade monetária — como poderia acontecer sob um padrão-ouro real. Mas dificilmente vale a pena fazê-lo simplesmente para se adaptar à média de quaisquer políticas que as autoridades monetárias do resto do mundo adotem. Muito melhor deixar o mercado, através de taxas de câmbio flutuantes, ajustar às condições mundiais os 5 por cento ou mais dos nossos recursos dedicados ao comércio internacional, enquanto reserva a política monetária para promover o uso eficaz dos 95 por cento.

Dos três guias listados, o nível de preços é claramente o mais importante por si só. Mantidas as demais condições, seria de longe a melhor das alternativas — como tantos economistas ilustres instaram no passado. Mas as outras coisas não são as mesmas. A ligação entre as ações de política da autoridade monetária e o nível de preços, embora inquestionavelmente presente, é mais indireta do que a ligação entre as ações de política da autoridade e qualquer um dos vários totais monetários. Além disso, a ação monetária leva mais tempo para afetar o nível de preços do que para afetar os totais monetários e tanto o atraso temporal quanto a magnitude do efeito variam com as circunstâncias. Como resultado, não podemos prever com precisão que efeito uma ação monetária específica terá sobre o nível de preços e, igualmente importante, quando terá esse efeito. Tentar controlar diretamente o nível de preços é, portanto, provável que torne a própria política monetária uma fonte de perturbação econômica devido a falsas paradas e partidas. Talvez, à medida que nosso entendimento dos fenômenos monetários avance, a situação mude. Mas no estágio atual do nosso entendimento, o caminho mais longo parece ser o caminho mais seguro para o nosso objetivo. Assim, acredito que um total monetário é o melhor guia ou critério imediato atualmente disponível para a política monetária — e acredito que importa muito menos qual total específico é escolhido do que que um seja escolhido.

Um segundo requisito para a política monetária é que a autoridade monetária evite oscilações bruscas na política. No passado, as autoridades monetárias ocasionalmente se moveram na direção errada — como no episódio da Grande Contração que enfatizei. Mais frequentemente, moveram-se na direção certa, embora muitas vezes tarde demais, mas erraram ao mover-se longe demais. Tarde demais e demais tem sido a prática geral. Por exemplo, no início de 1966, foi a política correta para o Federal Reserve mover-se em uma direção menos expansionista — embora devesse tê-lo feito pelo menos um ano antes. Mas quando se moveu, foi longe demais, produzindo a mudança mais acentuada na taxa de crescimento monetário da era do pós-guerra. Novamente, tendo ido longe demais, foi a política correta para o Fed inverter o curso no final de 1966. Mas novamente foi longe demais, não apenas restaurando, mas excedendo a taxa excessiva anterior de crescimento monetário. E este episódio não é exceção. Vez após vez este tem sido o curso seguido — como em 1919 e 1920, em 1937 e 1938, em 1953 e 1954, em 1959 e 1960.

A razão para a propensão a reagir exageradamente parece clara: a falha das autoridades monetárias em permitir o atraso entre as suas ações e os efeitos subsequentes na economia. Elas tendem a determinar as suas ações pelas condições de hoje — mas as suas ações afetarão a economia apenas seis ou nove ou doze ou quinze meses depois. Por isso sentem-se compelidas a pisar no freio, ou no acelerador, conforme o caso, com muita força.

Minha própria prescrição é ainda que a autoridade monetária vá até o fim para evitar tais oscilações, adotando publicamente a política de alcançar uma taxa constante de crescimento em um total monetário especificado. A taxa precisa de crescimento, como o total monetário preciso, é menos importante do que a adoção de alguma taxa declarada e conhecida. Eu mesmo argumentei por uma taxa que alcançaria em média a estabilidade aproximada no nível de preços dos produtos finais, que estimei exigiria algo como uma taxa de crescimento de 3 a 5 por cento ao ano em moeda corrente mais todos os depósitos em bancos comerciais ou uma taxa de crescimento ligeiramente inferior apenas em moeda corrente mais depósitos à vista. Mas seria melhor ter uma taxa fixa que produziria em média uma inflação moderada ou uma deflação moderada, desde que fosse estável, do que sofrer as amplas e erráticas perturbações que experimentamos.

Na falta da adoção de tal política publicamente declarada de uma taxa estável de crescimento monetário, constituiria uma melhoria importante se a autoridade monetária seguisse o ordenamento de auto-negação de evitar oscilações amplas. É um fato registrado que períodos de relativa estabilidade na taxa de crescimento monetário também foram períodos de relativa estabilidade na atividade econômica, tanto nos Estados Unidos como em outros países. Períodos de oscilações amplas na taxa de crescimento monetário também foram períodos de oscilações amplas na atividade econômica.

Ao estabelecer-se em um curso estável e manter-se nele, a autoridade monetária poderia fazer uma contribuição importante para promover a estabilidade econômica. Ao tornar esse curso um de crescimento constante, mas moderado, na quantidade de dinheiro, ela faria uma contribuição importante para evitar a inflação ou a deflação de preços. Outras forças ainda afetariam a economia, exigiriam mudança e ajuste e perturbariam o curso uniforme de nossos caminhos. Mas o crescimento monetário estável forneceria um clima monetário favorável à operação eficaz daquelas forças básicas de empreendedorismo, engenhosidade, invenção, trabalho árduo e parcimônia que são as verdadeiras fontes do crescimento econômico. Isso é o máximo que podemos pedir da política monetária no nosso estágio atual de conhecimento. Mas isso já é muito — e está claramente ao nosso alcance.

\section*{Referências}

\begin{enumerate}
    \item H. S. ELLIS, ed., \textit{A Survey of Contemporary Economics}. Philadelphia 1948.
    \item MILTON FRIEDMAN, "The Monetary Theory and Policy of Henry Simons," \textit{Jour. Law and Econ.}, Oct. 1967, 10, 1-13.
    \item MILTON FRIEDMAN, \textit{A Program for Monetary Stability}. New York 1959.
    \item E. A. GOLDENWEISER, "Postwar Problems and Policies," \textit{Fed. Res. Bull.}, Feb. 1945, 31, 112-21.
    \item P. T. HOMAN AND FRITZ MACHLUP, ed., \textit{Financing American Prosperity}. New York 1945.
    \item A. P. LERNER AND F. D. GRAHAM, ed., \textit{Planning and Paying for Full Employment}. Princeton 1946.
    \item J. S. MILL, \textit{Principles of Political Economy}, Bk. III, Ashley ed. New York 1929.
\end{enumerate}

\end{document}